\section{Experimental Objectives}

The experiments compare the direct image-space diffusion model with the wavelet-based diffusion model. The evaluation focuses on visual realism, convergence behavior, statistical similarity to validation data, preservation of vertical density profiles, and reproduction of scale-dependent fluctuation spectra.

The objective of the baseline diffusion model introduced in Chapter~3 is not to
reconstruct one specific DNS realization, but to generate new samples whose
distribution is close to that of the reference RTI dataset. Consequently, the
evaluation must be carried out at the level of \emph{ensembles} rather than on
individual images only.

In standard image generation tasks, this comparison is often performed with
feature-based distances such as the Fr\'echet Inception Distance (FID). Such
metrics are useful because they compare the statistics of two datasets in a
learned feature space. However, in the present context, the generated objects
are not natural images but scalar density fields arising from a physical
problem. It is therefore useful to complement generic perceptual metrics with
quantities that remain directly interpretable in terms of RTI structure.

For this reason, the evaluation adopted in this work combines two families of
metrics:

\begin{itemize}
    \item feature-based distribution distances, namely a classical
    Inception-based FID and a domain-adapted PhyFID;
    \item physically motivated one-dimensional summaries derived from the
    density fields, including line-wise mean profiles and fluctuation energy
    spectra.
\end{itemize}

This chapter presents the definition of these metrics and discusses the first
results obtained on the baseline SGM.

